\section{Related Work}

\subsection{Low-Rank Attention Compression}

\textbf{Linformer} (Wang et al., 2020) projects K and V to a fixed rank but requires 
retraining. \textbf{Nystr mformer} (Xiong et al., 2021) uses Nystr m approximation. 
\textbf{SOA Fused Attention} differs by factorizing the projections themselves 
rather than the attention matrix.

\textbf{ALBERT} (Lan et al., 2020) uses cross-layer parameter sharing, a form of 
low-rank structure. We use within-layer factorization.

\subsection{Weight Factorization}

\textbf{SVD-based compression} is well studied (Denton et al., 2014; Zhang et al., 
2015). We apply it to attention projections with a specific twist: the 
factorization is \textbf{zero-shot} (no calibration data needed) and combined 
with LoRA adapters for recovery.

\textbf{Tensor decomposition} (Novikov et al., 2015; Garipov et al., 2016) uses 
higher-order structure. Our approach uses standard matrix SVD.

\subsection{Parameter-Efficient Fine-Tuning}

\textbf{LoRA} (Hu et al., 2021) freezes base weights and trains low-rank adapters. 
We combine LoRA with SVD factorization: the base factors are SVD-initialized 
and frozen, while LoRA provides the recovery signal.

\textbf{AdaLoRA} (Zhang et al., 2023) adaptively allocates rank. \textbf{DoRA} 
(Liu et al., 2024) decomposes magnitude and direction. Our approach is a 
complementary technique: compress first, then apply LoRA.

\subsection{Attention Alternatives}

\textbf{Multi-Query Attention} (Shazeer, 2019) shares KV heads. **Grouped-Query 
Attention\textbf{ (Ainslie et al., 2023) groups KV heads. }Multi-Head Latent 
Attention** (DeepSeek-V2, 2024) projects KV to a low-rank latent. Our 
approach factorizes the projections themselves, which composes with GQA.

\subsection{Position of Our Work}

\begin{table}[H]
\centering
\small
\begin{tabular}{llllll}
\toprule
\textbf{Method} & \textbf{Calibration} & \textbf{Compresses} & \textbf{Zero-shot} & \textbf{GQA} & \textbf{RoPE} \\
\midrule
SVD (Vanilla) \& No \& All layers \& Yes \& Varies \& Varies \\
A3 \& Yes \& All layers \& No \& No \& No \\
AASVD \& Yes \& All layers \& No \& No \& No \\
LoRA \& No \& Nothing \& Yes \& Yes \& Yes \\
\textbf{SOA Fused} & \textbf{No} & \textbf{Q, K, V} & \textbf{Yes} & \textbf{Yes} & \textbf{Yes} \\
\bottomrule
\end{tabular}
\end{table}

\textbf{Key distinction.} SOA Fused Attention is the only method that combines:
\begin{itemize}
\item \textbf{Zero-shot} initialization (no calibration data)
\item \textbf{Real} parameter reduction (not just adapter)
\item \textbf{Attention-specific} targeting
\item \textbf{Full architecture support} (GQA, RoPE, causal)

\end{itemize}


\subsection{Comparison Table}

\begin{table}[H]
\centering
\small
\begin{tabular}{lllll}
\toprule
\textbf{Method} & \textbf{Reduction} & \textbf{Calibration} & \textbf{Retrain?} & \textbf{SST-2 acc} \\
\midrule
Baseline & 1$\times$ & -- & -- & 90.71\% \\
DistilBERT (distill) & 1.6$\times$ & -- \& Yes & 90.7\% \\
Vanilla SVD & 3$\times$ \& No \& LoRA & 84.5 \\
A3 & 3$\times$ \& Yes \& LoRA & 0.85 \\
\textbf{SOA Fused} & \textbf{3$\times$} & \textbf{No} & \textbf{LoRA} & \textbf{81.31} \\
\bottomrule
\end{tabular}
\end{table}

